\documentclass[11pt,a4paper]{article}

\usepackage[T1]{fontenc}
\usepackage[utf8]{inputenc}
\usepackage[english]{babel}
\usepackage{lmodern}
\usepackage[margin=2.4cm]{geometry}
\usepackage{microtype}
\usepackage{enumitem}
\usepackage{xcolor}
\usepackage{siunitx}
\usepackage{booktabs}
\usepackage{csquotes}
\usepackage[colorlinks=true,linkcolor=black,urlcolor=blue!55!black]{hyperref}
\usepackage{titlesec}

\definecolor{qcol}{HTML}{7A5C00}
\definecolor{vgood}{HTML}{1B6B2A}
\definecolor{vbad}{HTML}{A3231B}

\newcommand{\ann}[3]{%
  \medskip\noindent\textbf{[#1]}\quad p.~#2\quad\textcolor{qcol}{\textsf{\small #3}}\par\nopagebreak}
\newcommand{\note}[1]{\begingroup\small\color{qcol}\begin{quote}#1\end{quote}\endgroup}
\newcommand{\verdict}[1]{\textbf{\textcolor{vbad}{#1}}}
\newcommand{\ok}[1]{\textbf{\textcolor{vgood}{#1}}}

\titleformat{\section}{\large\bfseries}{\thesection.}{0.6em}{}
\setlength{\parskip}{0.35em}

\title{\textbf{Answers to the second review pass}\\[2pt]
  \large Questions from \texttt{main\_revision\_second.pdf} that are not settled inside the thesis}
\author{}
\date{2 August 2026}

\begin{document}
\maketitle
\vspace{-2.2em}

\noindent
Of the 294 annotations, 59 contain a question. Most were answered by clarifying the thesis
itself, and the row for each in \texttt{thesis/REVISION-log-second.md} says where. This
document holds what could not be settled that way: the two \emph{report-only} tasks whose
notes forbade editing before approval, and the questions that are decisions for you rather
than defects in the text.

\medskip\noindent
\textbf{Status.} The Step-1 reports below were produced first, as the notes require. Step~2
was then approved and \textbf{applied on 2 August 2026}; each report ends with what was
actually changed, and with the choices made where the evidence left more than one
defensible option.

\section{Annotation [36] --- Step-1 report on the glider-class table}

\note{\enquote{Task: fact-check a LaTeX table \dots\ Precision matters more than
completeness. \dots\ Do not edit anything before I approve Step 1.}}

\noindent
Step~2 has since been applied to Table~2.2, to its caption, and to annotation~[40], the
body sentence on the certification ladder.

\paragraph{Primary source used.} FAI Sporting Code, \emph{Common Section 7 --- Class O,
Hang Gliders and Paragliders}, \textbf{2024 Edition v3, effective 1 May 2024}, downloaded
from \texttt{fai.org}. Class definitions are in \S1.4.1 (pp.~11--12); the list of
subsections, including 7G, is in \S1.2. For the EN side, see item~6: the published
standard is a sold document, and what is cited instead is its publicly circulated final
draft.

\begin{enumerate}[leftmargin=1.4em, itemsep=0.5em]

\item \textbf{Is the numeric scale 1 / 1-2 / 2 / 2-3 an AFNOR classification?}
  \verdict{Refuted.} It is the \textbf{LTF} scale (the German
  \emph{Lufttüchtigkeitsforderungen}, successor to the DHV scale). The Sporting Code never
  writes \enquote{AFNOR}; it writes the pairs \enquote{EN~B/LTF~1--2} and
  \enquote{EN~C/LTF~2} when it defines the Class~3 sub-classes (\S1.4.1, p.~12). AFNOR was
  the older \emph{French} scheme and did not use those numbers.
  \textbf{The heading \enquote{EN/AFNOR ladder} was wrong and now reads EN/LTF.}

\item \textbf{Are the French parentheticals (\enquote{A ou 1}, \enquote{B ou 1-2}, \dots)
  documented?} \verdict{Not as written.} What is documented, in the Sporting Code itself,
  is the \emph{EN$\leftrightarrow$LTF} correspondence, not a French labelling. Presenting
  \enquote{A ou 1} as a classification label has no support in the primary source. They
  were removed.

\item \textbf{Does an official EN-to-numeric equivalence exist at all?}
  \textbf{Partially.} The Sporting Code asserts the equivalence only where it needs it, as
  the \emph{upper bound} of a sub-class (\enquote{Standard --- up to EN~B/LTF~1--2},
  \enquote{Sport --- up to EN~C/LTF~2}). That is a competition-eligibility statement, not a
  published one-to-one mapping of the two ladders, and the table no longer asserts it as an
  identity --- it records the two ceilings, which is exactly what the source states.

\item \textbf{Is \enquote{tandem / non-certified} one category or two?}
  \verdict{Two distinct things, wrongly merged.} They are not synonyms: a two-seater wing
  can be certified, whereas \enquote{non homologuée} means the wing carries no
  certification at all. The Sporting Code's Class~3 sub-classes are Open / Serial / Sport /
  Standard --- \emph{tandem does not appear}.

\item \textbf{\enquote{CCC --- CIVL Competition Class racing wings}.} \ok{Confirmed.}
  Sporting Code \S1.2: \enquote{Section 7G. CIVL Competition Class (CCC) --- Paragliders
  permitted in FAI Category~1 Cross-Country events (Class~3)}. The row is correct and was
  left untouched.

\item \textbf{\enquote{EN~A --- most stable, recreational}.} \ok{Substantially confirmed},
  on a second pass. The published EN~926-2:2013 is a sold, copyrighted CEN standard and I
  could not obtain it; but the \emph{final draft} submitted to formal vote,
  FprEN~926-2:2012, was circulated publicly by CEN/TC~136/WG~6, and its Table~1 gives the
  class descriptions verbatim: A is \enquote{maximum passive safety and extremely
  forgiving flying characteristics \dots\ designed for all pilots including pilots under
  all levels of training}. Your original wording was right; what was missing was a source
  for it. All four EN rows are now filled from that table, with the draft-versus-published
  caveat stated in the body rather than buried.

\item \textbf{Do \enquote{flexible wing} / \enquote{rigid wing} distinguish Class~1 from
  Class~2 from Class~5?} \verdict{Refuted --- this was the most serious error in the table.}
  The classes are defined by the \textbf{method of control}, and \emph{all three have a
  rigid primary structure} (\S1.4.1):
  \begin{itemize}[nosep, leftmargin=1.2em]
    \item \textbf{Class 1} --- rigid primary structure, \emph{pilot weight-shift as the
      sole method of control}; no pilot fairings or surrounding structures.
    \item \textbf{Class 2} --- rigid primary structure, \emph{movable aerodynamic surfaces
      as the primary method of control}.
    \item \textbf{Class 5} --- rigid primary structure, movable aerodynamic surfaces as the
      primary method of control \emph{in the roll axis}, plus the same prohibition on
      fairings and surrounding structures as Class~1.
  \end{itemize}
  So \enquote{Delta Class 1 / flexible wing} was wrong, and the identical note shared by
  Class~2 and Class~5 was not defensible: what separates them is the roll-axis restriction
  and the fairing prohibition. (Note also that \emph{Class~3 is paragliders} and
  \emph{Class~4} exists, for gliders that cannot take off or land in nil wind.)

\item \textbf{\enquote{Delta Class Sport}.} \ok{It is a class in its own right} ---
  explicitly \enquote{a sub-class of Class~1} (\S1.4.1, p.~11). But
  \verdict{\enquote{flexible wing, intermediate} was wrong}: the definition is not about
  performance. A Sport Class glider must meet the Class~1 definition \emph{and} carry an
  HGMA, BHPA or DHV type airworthiness certificate, be (or have been) commercially
  available for at least a year, be built of original parts only, and \emph{have a king
  post supporting the majority of the wing load when the wing is not flying}. It is a
  production-and-design restriction, not a performance tier.

\item \textbf{Does the FAI class scheme cover only hang gliders?} \verdict{No --- the
  heading was misleading.} The scheme is \textbf{Class~O}, which covers hang gliders
  \emph{and} paragliders; paragliders are \textbf{Class~3} inside the very same numbering.
  The heading now says so.
\end{enumerate}

\paragraph{Summary.} Of nine items: two confirmed (5, 8-partly), five refuted (1, 2, 4, 7,
9), one partial (3), and one (6) initially unverifiable but settled on a second pass
against the publicly circulated final draft of the standard.

\paragraph{Step 2, as applied.} Same booktabs style, two columns, row order unchanged, no
rows added. \emph{The one judgement call}: \enquote{tandem / non-certified} was kept as a
single row with the corrected note \enquote{two-seater and uncertified wings, pooled in
this field} --- splitting it would have added a row, which your scope rule forbids, and the
note is now true of the data even though the two things are distinct in law. The EN rows
were first left blank and then filled from Table~1 of the draft standard, once that source
turned up; the body now explains, in one paragraph, what a \enquote{sold} standard is and
why the draft is cited in place of the published text. Finally the caption: it claimed the
table was ordered \enquote{low to high performance}, which holds for EN/LTF and is false
for the FAI classes, so the caption now says so. The second column also became a wrapping
\texttt{p} column --- with real notes in it, the old \texttt{l} spec overflowed the text
block by \SI{135}{pt}.

\section{Annotation [66] --- Step-1 report on the multipath sentence}

\note{\enquote{Two elements are suspect. Verify each independently against the literature
before editing \dots\ Step~2, after my approval only.}}

\noindent
Text as it stood: \enquote{and \emph{multipath}---signals reaching the antenna only after
reflecting off terrain or the pilot's own equipment---adds high-frequency jitter of a few
metres.} Both of your suspicions were confirmed; Step~2 has since been applied.

\paragraph{Item 1 --- the parenthetical definition.} \verdict{The definition as written
describes NLOS, not multipath.} Multipath is the \emph{superposition} at the correlator of
the direct signal with delayed replicas; the direct signal is still present, and the error
comes from the distortion of the correlation function. The case in which the direct signal
is blocked and only a reflection is received is \emph{non-line-of-sight} (NLOS), a distinct
and separately treated phenomenon. The word \enquote{only} in the old sentence made it the
NLOS definition.

\paragraph{Item 2 --- \enquote{high-frequency jitter}.} \verdict{Refuted, and this is the
substantive error.} Code multipath is \emph{not} white and uncorrelated between epochs; it
is strongly time-correlated, and the current standards are criticised precisely for not
saying so. Garc\'ia Crespillo, Oliva Gonz\'alez and Caizzone (DLR), \emph{Airborne
Time-Correlated GNSS Multipath Error Modeling of Carrier-phase Smoothed Code}, ION ITM
2024, state in the Introduction that the avionics models (RTCA DO-253C, and the
elevation-dependent Gaussian models of Circiu et al.) \enquote{lack information about the
stochastic temporal behavior of the error, such as its time correlation}, and fit a
first-order Gauss--Markov process with a \textbf{time-correlation constant of
\SI{300}{\second}} to airborne data. The rate is set by the fringe frequency --- the rate
of change of excess path length divided by the carrier wavelength --- a
geometry-and-motion quantity, not a receiver one.

\emph{Caveats I will not paper over.} That study uses carrier-smoothed code on airliners
with aviation-grade antennas, not a free-flight variometer, and I could not settle two
things from primary sources: (i) the numerical fringe frequency for a receiver moving at
$\sim$\SI{10}{\metre\per\second} close to terrain, and hence whether \SI{1}{\hertz}
sampling can alias a fast-varying multipath error into something that \emph{looks} white;
(ii) whether \enquote{a few metres} is the right order for code multipath on consumer
hardware. Neither claim survives in the text, so nothing rests on them. What is safe to
say is that what dominates epoch-to-epoch scatter in \emph{open sky} is receiver tracking
noise scaled by the DOP, not multipath.

\paragraph{Replacement, as applied.}
\enquote{\dots and \emph{multipath}---the direct signal arriving together with replicas
delayed by reflection off terrain or the pilot's own equipment---biases the range on
timescales of minutes rather than adding independent noise from fix to fix.}
This keeps the em-dash apposition, corrects the definition, drops the unsupported jitter
claim, and states only what the source supports.

\paragraph{Also-check (reported, then acted on).} Two other sentences attributed the
observed positional scatter to multipath: the figure discussion in Sec.~2.6.1
(\enquote{large GNSS high-frequency noise, from multipath and the vertical dilution of
precision}) and the appendix remark on the hang-glider band. \textbf{Both were corrected in
the same pass}, since leaving them would have made the thesis contradict itself: a raised
\emph{floor} is by definition a white contribution, so it is receiver tracking noise scaled
by the DOP, not multipath. The appendix now also says plainly that why the hang-glider
fleet sits worse is \emph{not} settled by these spectra --- older or lower-grade receivers
being the obvious candidate --- rather than attributing it to multipath as before. No
white-noise offset is subtracted from the MSD at short lags anywhere in the text, so
nothing downstream depended on the wrong model.

\section{Questions that are decisions for you}

\ann{100}{17}{purple --- \enquote{Are we really sure we need this double check using also
the velocity bound?}}
Yes, and the thesis now argues it in place. Briefly: the Hampel test alone would delete
genuine circling fixes. At a \SI{10}{\second} cadence a thermalling circle of a few tens of
metres puts consecutive fixes on opposite sides of their own window median, so the residual
clears $\varepsilon_{\min}$ on geometry alone --- and those are exactly the trapping events
the transport analysis exists to measure. The speed bound is what separates that case from a
spike. Removing the corroboration would not simplify the rule; it would change what it
deletes, in the worst possible direction.

\ann{229}{29}{purple --- \enquote{siamo sicuri che un interpolatore lineare sia la scelta
migliore?}}
Per l'altitudine sì, ed è una scelta asimmetrica voluta (ora spiegata in
Sec.~2.7.6). La $z$ viene derivata una volta in più delle coordinate orizzontali prima di
essere usata --- $v_z$ è il discriminante della segmentazione --- quindi un artefatto di
interpolazione in $z$ finisce dentro l'assegnazione di fase, mentre lo stesso artefatto in
$E$ o $N$ viene assorbito dallo smoothing. L'interpolazione lineare non può fare overshoot
e ha un errore limitato da $\max|\ddot z|\,g^2/8$ (derivato ora in nota); una cubica
sarebbe più accurata dove $\ddot z$ è regolare e peggiore proprio dove non lo è, cioè alle
transizioni. Fra uno stimatore un po' conservativo ovunque e uno migliore in media ma che
può inventare una salita, per un canale che decide le fasi va scelto il primo.

\ann{239}{29}{red --- \enquote{molta attenzione a vedere come realmente è stata generata la
Figure 2.6}}
\verdict{Avevi ragione: il difetto c'era.} La figura marcava il taglio a
$10\,\Delta t$, cioè la sola parte relativa del criterio, mentre il bound effettivo è
$g_{\max}=\min(10\Delta t,\ \max(\SI{20}{\second},2\Delta t))$. Le due coincidono
\emph{solo} alla cadenza di \SI{1}{\second}: da \SI{2}{\second} a \SI{10}{\second} vincola
il cap assoluto, oltre il floor $2\Delta t$. La linea tratteggiata quindi non rappresentava
il criterio in vigore per la maggior parte dei voli. Corretto nel codice e rigenerata: il
pannello~(a) ora normalizza il gap di ogni volo al \emph{suo} $g_{\max}$, così il criterio
vero cade esattamente a $1$.

\ann{76}{14}{orange --- \enquote{Are we 100\% sure of these numbers?}}
I numeri sono macro generate da \texttt{generate\_census\_stats.py} a partire dallo scan
completo dell'archivio, quindi non sono trascritti a mano e non possono divergere dal dato.
Il limite vero è un altro, ed è ora dichiarato in appendice: lo scan gira su tracce
\emph{parsate}, non ancora pulite, quindi ogni conteggio del Cap.~2 va riletto dopo il
fix-level cleaning. Fino ad allora vanno considerati stime correnti, non cifre definitive.

\ann{237}{29}{yellow --- \enquote{È il posto giusto dove mettere questo paragrafo?}}
Sì, l'ho lasciato dov'è. Il paragrafo definisce $g_{\max}$ nel punto in cui il bound viene
usato per la prima volta, ed è anche il punto in cui serve al lettore per capire cosa
distingue un buco interpolato da uno split. Spostarlo più avanti costringerebbe ad
anticipare il simbolo; spostarlo prima lo separerebbe dalla regola che lo usa.

\ann{265}{32}{yellow --- la mappa dei siti di decollo}
Specificata in Sec.~2.8, ma non ancora prodotta: richiede codice nuovo, non solo una
rigenerazione. Una nota di merito: la mappa non può essere \enquote{della Francia} e basta,
perché la CFD ammette voli fatti all'estero e le sorgenti future (WeGlide, XContest) non
sono francesi. Il testo la definisce quindi sul bounding box effettivo delle coordinate di
decollo, con la Francia disegnata dentro.

\section{What was changed in code, not only in text}

\begin{itemize}[leftmargin=1.3em, itemsep=0.3em]
  \item \texttt{BARO\_PRESENT\_MIN} raised \num{0.5} $\to$ \num{0.95} and made a single
    source of truth in \texttt{soaring.analysis.altitude\_noise}, imported by
    \texttt{preprocessing} so cleaning, census and PSD cannot drift; tests updated ([293]).
    \emph{Measured effect}: it reclassified \num{241} paraglider flights and no
    hang-glider ones, \SI{0.13}{\percent} of the archive --- which is the direct
    confirmation of the bimodality the appendix asserts. It also cut
    \texttt{ParaBaroMissMaxFixes} from \num{14878} to \num{1397}, i.e.\ most of what the
    text used to call a \enquote{rare pathological tail} was a classification artefact of
    the old threshold, not a sensor one. Both sentences describing it were rewritten.
  \item \texttt{make\_flightlevel\_diagnostics\_figure}: 2$\times$2 $\to$ 2$\times$3 panels,
    adding the altitude-range distribution and its retention curve. No rescan needed --- the
    range is already in the cached census ([181], [185]).
  \item \texttt{make\_gap\_diagnostics\_figure}: panel~(a) now normalised by each flight's
    own $g_{\max}$ ([239], above).
  \item \texttt{generate\_census\_stats.py}: new \texttt{\textbackslash StatCat*} family
    reading \texttt{catalog.csv}, for the placeholder-date census of [50]. \emph{Result}:
    \num{8} of \num{203343} paraglider entries carry the \texttt{0000-00-00} sentinel
    (\SI{0.0039}{\percent}) and none of the hang-glider ones; the worst-affected season
    loses \SI{0.17}{\percent} of its entries. Your worry that one early season might be
    gutted by the cut is empirically unfounded, and the text now says so with the numbers.
\end{itemize}

\section{Two things I chose, that you may want to overrule}

\begin{enumerate}[leftmargin=1.4em, itemsep=0.4em]
  \item \textbf{\texttt{BARO\_PRESENT\_MIN} is \texttt{>= 0.95}, not \texttt{> 0.95}.} Your
    note said \enquote{sopra il 95\,\%}. A flight covering exactly \SI{95}{\percent} is
    admitted. This matches how every other threshold in the pipeline is written (inclusive)
    and, given that presence is bimodal, no flight is expected to sit exactly on the
    boundary anyway. One character to change if you disagree.
  \item \textbf{Pass-1 markup that carried no green annotation was left coloured.} Black now
    means \enquote{accepted in this pass}. The alternative reading --- that anything you did
    not object to is accepted --- would have turned most of the chapter black and lost the
    distinction. Say so and I flip it.
\end{enumerate}

\end{document}
